%  OPJS 1763 September 15, 2025.
% subject class and Key words??
\documentclass[reqno]{amsart}
\usepackage{a4wide}
\usepackage{hyperref}

\AtBeginDocument{{\noindent\small
\emph{Electronic Journal of Differential Equations},
Vol. 2026 (2026), No. 01, pp. 1--13.\newline
ISSN: 1072-6691. URL: https://ejde.math.txstate.edu, DOI: 10.58997/ejde.2026.01}
\thanks{\copyright 2026. This work is licensed under a CC BY 4.0 license.}
\vspace{8mm}}

\begin{document}
\title[\hfilneg EJDE-2026/01\hfil $n$-dimensional stationary Navier-Stokes equations]
{Variational approach for the $n$-dimensional stationary Navier-Stokes equations
with a damping term}

\author[A. Khatib, A. Moameni, S. Mousavinasr \hfil EJDE-2026/01\hfilneg]
{Alireza Khatib, Abbas Moameni, Somayeh Mousavinasr}

\address{Alireza Khatib \newline
Universidade Federal do Amazonas, Manaus, AM, Brazil}
\email{alireza@ufam.edu.br}

\address{Abbas Moameni \newline
School of Mathematics and Statistics,
Carleton University, Ottawa, ON,  Canada}
\email{momeni@math.carleton.ca}

\address{Somayeh Mousavinasr \newline
Universidade Federal do Amazonas,
Manaus,  AM, Brazil}
\email{somayeh@ufam.edu.br}

\thanks{Submitted September 15, 2025. Published January 6, 2026.}
\subjclass[2020]{35Q30, 35A15}  
\keywords{Navier-Stokes equation; variational method; damping term}

\begin{abstract}harvesting effort
We study the $n$-dimensional stationary Navier-Stokes equations with a damping
term by developing a new general minimax principle. This principle is sufficiently
broad to be applied in various contexts, and here it is used to establish the existence
of weak solutions for both linear and nonlinear damping, without restrictions on the
damping constant. The damping term, which models physical effects such as porous media
flow, drag, friction, and dissipation, also provides a mathematical advantage by
improving the regularity of solutions compared to the classical Navier-Stokes system.

Our results cover the cases of positive and negative damping constants and
yield existence theorems under different ranges of $p$ and spatial dimensions.
In particular, we prove solvability even in borderline situations,
such as when $\mu = -\lambda_1$, where coercivity is lost and traditional minimax
arguments typically fail. The general minimax framework we introduce is flexible and
can be adapted to other nonlinear PDEs, especially when symmetry or structural properties
are involved.
\end{abstract}

\maketitle
\numberwithin{equation}{section}
\newtheorem{theorem}{Theorem}[section]
\newtheorem{lemma}[theorem]{Lemma}
\newtheorem{corollary}[theorem]{Corollary}
\newtheorem{remark}[theorem]{Remark}
\allowdisplaybreaks

\section{Introduction}

In this article we study  the $n$-dimensional  stationary Navier-Stokes equations
with a damping term,
\begin{equation}\label{eq1}
 \begin{gathered}
 -\Delta u+ (u \cdot \nabla)u  +\mu |u|^{p-2}u = f(x) + \nabla P  \quad
   \forall x\in \Omega\\
 \nabla  u=0  \quad  \forall x\in \Omega \\
  u=0  \quad   \forall x\in \partial\Omega
\end{gathered}
\end{equation}
where \(\Omega \subset \mathbb{R} ^n\) is  bounded,  $p\geq 1$ and
$\mu \in \mathbb{R}$.  We address both  linear and nonlinear dampings
and we are allowing $\mu$ to take both positive and negative values.
Here \(u  = (u_1, u_2,\dots, u_n)\) is the velocity, \(P\) stand for scalar pressure
and \(f\) is the  external force.

The analysis of the Navier-Stokes equations is a central theme in mathematical fluid
mechanics. For the classical evolutionary system without damping, the existence of
global weak solutions was established by Leray \cite{JL} and Hopf \cite{EH}, but the
uniqueness of weak solutions and the global existence of strong solutions remain open
problems. These longstanding challenges have motivated researchers to study modified
models where additional terms improve the mathematical structure of the equations.

One such modification is the inclusion of a damping term of the form $\mu |u|^{p-2}u$.
From the physical viewpoint, this term models resistance to motion and arises naturally
in contexts such as porous media flows, drag or friction effects, and other dissipative
mechanisms (see \cite{LHS, RTE}). From the mathematical viewpoint, damping often
leads to better control of solutions, sometimes yielding results that are out of reach
for the standard Navier-Stokes equations. This has stimulated extensive research on the
evolutionary Navier-Stokes system with damping, leading to results on global existence,
regularity, decay rates, attractors, and stability (see, e.g.,
\cite{XC, YJX, ZHJ, ZHM, XLS,ZJZ, YZ, ZZC}).

In contrast,  stationary Navier-Stokes equations with damping have received comparatively
less attention. Some existence and uniqueness results are available for the case
$\mu > 0$ \cite{WLW}, and there is growing literature on numerical approaches
(see \cite{ZLD, LML, HQL, HQY, ZBS}). However, general variational approaches capable
of handling both positive and negative damping, as well as borderline cases where
coercivity is lost, are still lacking.

In this work, we develop a new minimax principle on convex subsets of Banach spaces and
apply it to the $n$-dimensional stationary Navier-Stokes equations with damping.
Our framework is broad and flexible: depending on the choice of convex set, it allows
us to obtain solutions with additional structural properties, such as symmetry.
With this method we prove existence results for both linear ($p=2$) and nonlinear
damping, without restriction on the damping constant, and even in the critical case
$\mu = -\lambda_1$. This general minimax principle is of independent interest and may
find applications to other nonlinear PDEs beyond the Navier-Stokes system.

In this work,  we  consider the Banach space
\[
V= \{ u\in H_0^1(\Omega)\cap L^p(\Omega),\nabla . u=0\},
\]
equipped with the  norm
\[
\rVert u \rVert:= \rVert u \rVert_{ H_0^{1}(\Omega)}
+\rVert u \rVert_{ L^{p}(\Omega)}.
\]
Let \(\Lambda u\) be the operator \(\Lambda u := (u . \nabla)u\),
and \( K\) be a convex and weakly closed subset of \(V\).
We shall define \( M : K \times K \to\mathbb{R}\) as follows,
\begin{equation}\label{def II}
    M(u,v)=\frac{1}{2}\int_{\Omega} |\nabla u|^2\,dx
    -\frac{1}{2}\int_{\Omega} |\nabla v|^2\,dx
    +\int_\Omega (\Lambda u-f(x)  -\frac{1}{p}|u|^{p-2}u) (u-v)\,dx,
\end{equation}
where \(f \in L^2(\Omega)\). The following variational principle on general convex
sets $K$  is a  key component in our arguments. It is also broad enough to deal with
various other cases by choosing a convex set $K$ accordingly.

\begin{theorem}\label{variational}
Let  \( K\) be a convex and weakly closed subset of \(V\). Assume that the following
two assertions hold:
\begin{itemize}
\item[(i)] There  exists \(\bar{u} \in K\) for   such that
    \[
    M(\bar{u}, v)\leq 0,\quad  \forall v\in  K,
    \]
    where \(M \) is defined in \eqref{def II}.

\item[(ii)] There exists \(\bar{v} \in K\) such that \[-\Delta \bar{v} + \nabla P = f(x) +|\bar{u}|^{p-2}\bar{u} - \Lambda \bar{u},\]
    in the weak sense, i.e.,
\[
   \int_\Omega\nabla\bar{v}.\nabla\eta\,dx
   =\int_\Omega (f(x)  +|\bar{u}|^{p-2}\bar{u} -
   \Lambda \bar{u})\eta\,dx,\quad \forall\eta\in V.
\]
\end{itemize}
Then \(\bar{u} \in K\) is a weak solution of the equation
\[
-\Delta u+ \;\Lambda u \; = f(x) + \nabla P +|u|^{p-2}u \,.
\]
\end{theorem}
It is worth noting  that the primary consequence of this theorem centres on the choice of
$K$, i.e., by choosing an appropriate $K$, one is able to establish the existence of a
solution enjoying all the
properties induced by the set $K$ (see Remark \ref{symm} for an application where the
problem \eqref{eq1} has some symmetry properties).
Also,  Condition (i) in Theorem \ref{variational} is most of the time guarantied
because of the well-known Ky Fan's min-max principle by Brezis-Nirenberg-Stampacchia
\cite{HB}.
We provide more details of how to apply the above theorem in the sequel.
As an application of the above Theorem we first prove the following result.

\begin{theorem}\label{main theorem 1}
Let \(\Omega\) be a  bounded \(C^2\) domain in \(\mathbb{R} ^n\) and  \(\mu<0\).
Then for \(f\in L^2(\Omega) \) small enough,  the following statements hold:
 \begin{itemize}
 \item[(i)] For \(n\leq 4\) and \(p>2\), the Navier-Stokes equation
     \eqref{eq1} has a  solution \(u \in  W^{2,2}(\Omega)  \).

 \item[(ii)] For \(5\leq n\leq 7\) and  \( 2< p \leq \frac{2n-4}{n-4}\),
 the Navier-Stokes equation \eqref{eq1} has a  solution in \( W^{2,2}(\Omega)\).
 \end{itemize}
In both cases there exists a scalar function $P:\Omega \to \mathbb{R}$ and a constant $C>0$
such that
\begin{equation} \label{nabla11}
     \rVert \Delta {u} \rVert_{L^2(\Omega)}+\rVert \nabla P \rVert_{L^2(\Omega)}
     \leq C\big( \rVert f\rVert_{L^2(\Omega)}+\rVert u\rVert^{p-1}_{L^{2(p-1)}}
     +\rVert {u}\rVert^{2}_{W^{1,2^*}(\Omega)} \rVert  {u}\rVert^2_{L^{\frac{n}{2}}(\Omega)} \big) .
\end{equation}
  where $2^*=2n/(n-2)$.
\end{theorem}

When the constant $\mu$ in the damping term  is non-negative  we can cover
higher values for $p$  as shown in the following theorem.

\begin{theorem} \label{main theorem 2}
Let \(\Omega \subset \mathbb{R}^n\) be  a smooth bounded domain and \(\mu>0\).
Suppose that \( p\geq1\)  and \(f\in L^2(\Omega)\). Then there exists \(u \in V\) such
that the following holds:
 \begin{itemize}
\item[(i)] If \(n\geq 2\), then
 \[ \int_\Omega\nabla u.\nabla\eta\,dx
 + \int_\Omega |{u}|^{p-2}u\; \eta  + \int_\Omega \Lambda u\;\eta\,dx
 =\int_\Omega f(x)\eta\,dx,\quad \forall\eta\in  C_c^1(\Omega),  \textit{ with }
  \nabla. \eta=0.
 \]

\item[(ii)] If \(n\leq 4 \) or \(p\geq 4\), then
\[
\int_\Omega\nabla u.\nabla\eta\,dx
+ \int_\Omega |{u}|^{p-2}u\; \eta \,dx+ \int_\Omega \Lambda u\;\eta\,dx
=\int_\Omega f(x)\eta\,dx,\quad \forall\eta \in V.
\]
\end{itemize}
\end{theorem}

We would like to remark that the solution we are getting in part (i) of the above
theorem is weaker than the one we are getting in part (ii).
This is so because all the test functions  $\eta$ in  part $(i)$ are coming from
$C_c^1(\Omega)$ on contrary to part $(ii)$  where the test functions  $\eta$
live in a less regular space $ H_0^1(\Omega)\cap L^p(\Omega)$.

We shall also deal with the linear damping term where $p=2$ for positive and negative
values of $\mu$.
To state our result we first recall the following standard fact about the first
eigenfunction  of the Laplacian on bounded domains. Recall that
\[
 \lambda_1 = \min_{\psi\in H_0^1(\Omega)\backslash \{0\}}
 \frac{\int_{\Omega} |\nabla \psi|^2\,dx}{\int_\Omega  |\psi|^{2}\,dx}.
\]
where the minimum is taken over all $\psi: H_0^1(\Omega)\to\mathbb{R} $.
Note that in Theorem \ref{main theorem 2} we have already covered the case $\mu>0$.
Here is our result for the linear case where we are allowing negative values for $\mu$.

\begin{theorem}\label{main theorem 3}
Let \(\Omega\) be  smooth bounded domain in \(\mathbb{R} ^n\)  and  \(p=2\).
Assume that \(-\lambda_1\leq\mu<0\),  and \(f\in L^2(\Omega) \).
Then there exists \(u \in V\) such that the following  assertions hold:
\begin{itemize}
\item[(i)] If \(n\geq 2\), then  \[\int_\Omega\nabla u.\nabla\eta\,dx
+\mu \int_\Omega u\eta\,dx  + \int_\Omega \Lambda u\;\eta\,dx
=\int_\Omega f(x)\eta\,dx,\quad \forall\eta\in  C_c^1(\Omega),\quad
 \textit{with } \nabla . \eta=0.
 \]

\item[(ii)] If \(n\leq 4 \), then \[\int_\Omega\nabla u.\nabla\eta\,dx
+ \mu\int_\Omega u\; \eta\,dx + \int_\Omega \Lambda u\;\eta\,dx
 =\int_\Omega f(x)\eta\,dx,\quad \forall\eta \in V.
\]
\end{itemize}
\end{theorem}

The highlight of the above theorem is the case where $\mu=-\lambda_1$ in which case
one losses the coercivity required in most minimax arguments.

\begin{remark}\label{symm} \rm
Even though our main objective in this paper is to prove existence results having
a damping term in mind,  we would like emphasize that the applications of Theorem
\ref{variational}  goes well beyond this goal.  In light of this remark, let us
define the maps $\pi_1,\pi_2,\pi_3: \Omega \subset \mathbb{R}^3 \to\Omega$ as follow:
\begin{gather*}
    \pi_1(x_1,x_2,x_3) = (-x_1,x_2,x_3), \\
    \pi_2(x_1,x_2,x_3) = (x_1,-x_2,x_3), \\
    \pi_3(x_1,x_2,x_3) = (x_1,x_2,-x_3).
\end{gather*}

We consider the 3D case of the stationary Navier-Stokes equations with damping
presented in equation \eqref{eq1}. Assume that $\Omega$ is invariant under the maps
$\pi_1, \pi_2, \pi_3: \Omega \to\Omega$.
Moreover, assume that $K_S$ is a subset of $V$ containing all $ u \in V$ with the
following properties:
\begin{equation}\label{eq:relation1}
\begin{gathered}
u_1(x_1,x_2,x_3) = -u_1(-x_1,x_2,x_3), \\
u_2(x_1,x_2,x_3) = u_2(-x_1,x_2,x_3), \\
u_3(x_1,x_2,x_3) = u_3(-x_1,x_2,x_3).
\end{gathered}
\end{equation}
Furthermore, assume that $f(x)\in L^2(\Omega)$ also holds the same properties; i.e.,
\begin{gather*}
f_1(x_1,x_2,x_3) = -f_1(-x_1,x_2,x_3), \\
f_2(x_1,x_2,x_3) = f_2(-x_1,x_2,x_3), \\
f_3(x_1,x_2,x_3) = f_3(-x_1,x_2,x_3).
\end{gather*}
Then, the solution $u=(u_1,u_2,u_3) $ obtained in Theorems \ref{main theorem 1},
\ref{main theorem 2} and \ref{main theorem 3}  is symmetric in the sense
\eqref{eq:relation1}. Indeed, the symmetry of solutions follows from the uniqueness
of the solution to the corresponding linear problem.
\end{remark}

The article is organized as follows. In section 2,
we  prove  Theorems \ref{variational} and  \ref{main theorem 1}
through a minimax principle.
Section 3 is devoted to the proof of our results in Theorems
\ref{main theorem 2} and \ref{main theorem 3}.


\section{A minimax principle and the proof of Theorem \ref{main theorem 1}}

In this section, we first prove an adapted version of  variational principle
presented in Theorem \ref{variational} which is  applicable specifically to our
problem when \(\mu<0\), and $p>2$. Afterwards, we proceed with the proof of
Theorem \ref{main theorem 1}.

 We  consider the Banach space
\(
V= \{ u\in H_0^1(\Omega)\cap L^p(\Omega),\nabla . u=0\}
\)
equipped with the  norm
\[
\rVert u \rVert:= \rVert u \rVert_{ H_0^{1}(\Omega)}
+\rVert u \rVert_{ L^{p}(\Omega)}.
\]

Let \(\Lambda u\) be the operator \(\Lambda u := (u . \nabla)u\), that is
\[
\langle \Lambda u, v\rangle
= \int_\Omega (\Lambda u)v =\int_\Omega\sum_{j,k=1}^n
u_k\frac{\partial u_j}{\partial x_k}v_j.
\]

Let \( K\) be a convex and weakly closed subset of \(V\).
As stated in Theorem \ref{variational} we shall consider the functional
\( M : K \times K \to\mathbb{R}\) given in  \eqref{def II}.

\begin{proof}[Proof of Theorem \ref{variational}]
It follows from condition (i) in the theorem that there exists \(\bar{u} \in K\)
such that
\begin{equation}\label{ineq 1}
    \frac{1}{2}\int_{\Omega} |\nabla \bar{u}|^2\,dx
    -\frac{1}{2}\int_{\Omega} |\nabla v|^2\,dx
    \leq \int_\Omega (   f(x)+\frac{1}{p}|\bar{u}|^{p}- \Lambda \bar{u}) (\bar{u}-v)\,dx ,
    \quad \forall v\in  K.
\end{equation}
It also follows from (ii) that there exists \(\bar{v} \in K\) such that
\begin{equation}\label{ineq 2}
\int_\Omega\nabla\bar{v}.\nabla\eta\,dx=\int_\Omega (f(x)
+|\bar{u}|^{p-2}\bar{u} - \Lambda \bar{u})\eta\,dx,\quad \forall\eta\in V.
\end{equation}
Substituting \(\eta = \bar{u} - \bar{v}\) in the latter equality gives
\begin{equation}\label{ineq 3}
\int_\Omega\nabla\bar{v}.\nabla(\bar{u} - \bar{v})\,dx
=\int_\Omega (f(x)  +|\bar{u}|^{p-2}\bar{u} - \Lambda \bar{u})(\bar{u}
- \bar{v})\,dx,\quad \forall\eta\in V.
\end{equation}
Setting \(v = \bar{v}\) in \eqref{ineq 1} and taking into account the equality
\eqref{ineq 3} we obtain that
\begin{equation}\label{ineq 4}
 \int_\Omega\nabla\bar{v}.\nabla(\bar{u} - \bar{v})\,dx
 \geq\frac{1}{2}\int_{\Omega} |\nabla \bar{u}|^2\,dx
 -\frac{1}{2}\int_{\Omega} |\nabla \bar{v}|^2\,dx.
\end{equation}
On the other hand, it follows from the convexity of \( g(t)=\frac{1}{2}t^2\) that
\begin{equation}\label{ineq 5}
    \int_\Omega\nabla\bar{v}.\nabla(\bar{u} - \bar{v})\,dx \leq\frac{1}{2}\int_{\Omega} |\nabla \bar{u}|^2\,dx -\frac{1}{2}\int_{\Omega} |\nabla \bar{v}|^2\,dx.
\end{equation}
Inequalities \eqref{ineq 4} and \eqref{ineq 5} together imply that
\[
\int_\Omega\nabla\bar{v}.\nabla(\bar{u} - \bar{v})\,dx
=\frac{1}{2}\int_{\Omega} |\nabla \bar{u}|^2\,dx-\frac{1}{2}\int_{\Omega}
|\nabla \bar{v}|^2\,dx.
\]
Therefore,
\[
\int_{\Omega} |\nabla \bar{u}-\nabla \bar{v}|^2=0,
\]
from which it follows that \(\bar{v}=\bar{u}\) for a.e.\
 \(x \in\Omega\). Hence, the equality \eqref{ineq 2} proves the desired  result.
\end{proof}

We shall apply Theorem \ref{variational} to prove the existence of solution  in
Theorem \ref{main theorem 1}.  The convex subset \(K\) of \(V\)
required in Theorem \ref{variational} is defined by
\begin{equation}\label{def K}
K(r)=\{u\in V: \rVert u \rVert_{ W^{2,2}(\Omega)}\leq r\},
\end{equation}
for some \(r > 0\) to be determined. To see that \(K(r)\) is weakly closed,
we present the proof of this statement in the following lemma.

\begin{lemma}\label{Kclosed}
Let \(r > 0\) be fixed. The set
\[
K(r)=\{u\in V: \rVert u \rVert_{ W^{2,2}(\Omega)}\leq r\}
\]
is weakly closed in \(V\).
\end{lemma}

\begin{proof}
Let \(\{u_m\}\) be a sequence in \(K(r)\) such that \(u_m\rightharpoonup u\) weakly
in \(V\).
Then there exists a subsequence of \(u_m\), denoted by \(u_m\) again such that
\(u_m\to u\) a.e in \(\Omega\).
On the other hand,  \(\rVert u_m \rVert_{ W^{2,2}(\Omega)}\leq r\) for all
\(m \in \mathbb{N}\) and so  \(\{u_m\}\)
is bounded in \(W^{2,2}(\Omega)\). Going if necessary to a subsequence, there exists
\(\bar{u}  \in W^{2,2}(\Omega)\)  such that \(u_m \rightharpoonup \bar{u}\) weakly in
\(W^{2,2}(\Omega)\) and \(u_m(x) \to\bar{u}(x) \) for a.e.\
 \(x \in \Omega\). It follows then \(u(x)=\bar{u}(x) \) for a.e. \(x \in \Omega\).
Thus \(u_m \rightharpoonup u\) weakly in \(W^{2,2}(\Omega)\). Now from the
weak lower semi-continuity of the norm in \(W^{2,2}(\Omega)\)  follows that
\[
\rVert u \rVert_{ W^{2,2}(\Omega)}
\leq \liminf_{m\to\infty} \rVert u_m \rVert_{ W^{2,2}(\Omega)}\leq r ,
\]
which means that \(u\in K(r)\).
\end{proof}

To apply Theorem \ref{variational}, we need to verify both conditions (i)
and (ii) in this Theorem. To verify
condition (i) we  use the following version of the well-known Ky Fan's
min-max principle \cite{HB}. We refer to \cite[Lemma 12.1]{NG} for a proof.

\begin{lemma}\label{Ky Fan 2}
Let \(E\) be a closed convex subset of a reflexive Banach space \(H\), and consider
\(M : E \times E \to\bar{\mathbb{R}}\)
to be a functional such that:
\begin{itemize}
    \item[(1)] For each \(y\in E\), the map \(x \to M(x, y)\) is weakly lower semi-
     continuous on \(E\).
    \item[(2)] For each \(x \in E\), the map \(y \to M(x, y)\) is concave on \(E\).
    \item[(3)] There exists \( \gamma\in \mathbb{R}\) such that \(M(x, x) \leq \gamma\)
    for every  \(x \in E\).
    \item[(4)] There exists a \(y_0 \in E\) such that
     \(E_0 = \{x \in E : M(x, y_0) \leq \gamma\} \) is bounded.
\end{itemize}
Then there exits \(\bar{x} \in E\) such that \(M(\bar{x}, y) \leq \gamma\) for all
\(y \in E\).
\end{lemma}

One of the requirements in Lemma \ref{Ky Fan 2} is  the lower semi-continuity of
$M(.\;,v)$ for a fixed $v$. To verify  that, we begin with the following Lemma.

\begin{lemma}\label{lem SC 1}
For each \( v\in K\), the map   \( u\to\langle \Lambda u, v\rangle  \) is weakly
continuous on \(K\) for the values of  \( n,p \) in Theorem \ref{main theorem 1}.
\end{lemma}

\begin{proof}
Fix \( v \in  K\), and let  \(u^m\rightharpoonup u\) weakly in \(K\). We have
\begin{align*}
\big| \langle \Lambda u^m, v\rangle -  \langle \Lambda u, v\rangle  \big|
&=\Big|\sum_{j,k=1}^n\int_\Omega\big(u_k^m\frac{\partial u_j^m}{\partial x_k}v_j
 -u_k\frac{\partial u_j}{\partial x_k}v_j\big) \,dx\Big| \\
&=\Big|\sum_{j,k=1}^n\int_\Omega\big( (u_k^m-u_k)\frac{\partial u_j^m}{\partial x_k}v_j
 +u_k\frac{\partial (u_j^m-u_j)}{\partial x_k}v_j\big) \,dx\Big| \\
&\leq \sum_{j,k=1}^n\int_\Omega\Big| (u_k^m-u_k)\frac{\partial u_j^m}{\partial x_k}v_j
 \Big|
+\sum_{j,k=1}^n\int_\Omega \Big| u_k\frac{\partial (u_j^m-u_j)}{\partial x_k}v_j\Big|
\,dx.
\end{align*}
On the other hand, by H\"{o}lder inequality we conclude that
\[
\sum_{j,k=1}^n\int_\Omega\Big| (u_k^m-u_k)\frac{\partial u_j^m}{\partial x_k}v_j\Big|
    \leq \rVert (u^m-u)v \rVert_{L^2}\rVert \nabla u^m \rVert_{L^2}
    \leq \rVert u^m-u \rVert_{L^4}\rVert v \rVert_{L^4} \rVert \nabla u^m \rVert_{L^2}.
\]
Therefore,
\[
\big| \langle\Lambda u^m, v\rangle -  \langle\Lambda u, v\rangle \big|
\leq \rVert u^m-u \rVert_{L^4}\rVert v \rVert_{L^4} \rVert \nabla u^m \rVert_{L^2}
 +\sum_{j,k=1}^n\int_\Omega \Big| u_k\frac{\partial (u_j^m-u_j)}{\partial x_k}v_j\Big|
 \,dx.
\]
Moreover, since the space \(W^{2,2}(\Omega)\) is compactly imbedded into
\(L^4(\Omega)\), for all \(n\leq 7\),  it follows that \(u^m\to u\) strongly
in \(L^4(\Omega)\).
Furthermore, since $u,v$ are in \(W^{2,2}(\Omega)\), we deduce
from H\"{o}lder's inequality that \(u_kv_j\in L^2(\Omega)\).
Finally,  since  \(\nabla u^m\rightharpoonup \nabla u\) weakly in \(K\),  by
definition of weak convergence  the result follows.
\end{proof}

\begin{lemma} \label{WLSC}
Fro each  \( v\in K\), the map \(u\to M(u, v) \) is weakly lower semi-continuous on
\(K\) for the values of  \( n,p \) in Theorem \ref{main theorem 1}.
\end{lemma}

\begin{proof}
 Let \( v \in  K\) be fixed and \(u_m\rightharpoonup u\) weakly in \(K\).
Since $<\Lambda u , u> = 0 $ resulting from \eqref{def II} we have
\begin{equation}\label{M(u)}
\begin{aligned}
M(u,v)
& =\frac{1}{2}\int_{\Omega} |\nabla u|^2\,dx -\langle\Lambda u , v\rangle
 -\int_{\Omega} f(x)u\,dx-\frac{1}{p}\int_{\Omega}  |u|^{p}\,dx  \\
&\quad +  \frac{1}{p}\int_{\Omega}  |u|^{p-2}uv\,dx+ \frac{1}{2}
\int_{\Omega} |\nabla v|^2\,dx+\int_{\Omega} f(x)v\,dx,
\end{aligned}
\end{equation}
Now we shall verify lower semi-continuity of every single part in \eqref{M(u)} separately.
Note that the last two terms in \eqref{M(u)} are constant  with respect to  $u$.
\begin{itemize}
 \item Since the function \(g(u) = |u|^2 \) is convex, it can easily be shown that
\[
 \int_{\Omega}|\nabla u|^2\,dx\leq \liminf_{m \to \infty}\int_{\Omega}|\nabla u^m|^2\,dx .
\]
 that implies the map \( u\to\int_{\Omega}|\nabla u|^2  \,dx\) is weakly lower
 semi-continuous.

 \item The map \( u\to-\int_{\Omega}\Lambda u. v \,dx\) is weakly lower semi-continuous
 by Lemma \ref{lem SC 1}.

 \item  Since \(f \in L^2(\Omega)\), applying the definition of weak convergence
 leads to
 \[
 \int_{\Omega}f(x)\  u  \,dx=\liminf_{n\to \infty}\int_{\Omega}f(x) u^m\,dx.
 \]

 \item The map \( u\to\int_{\Omega}|u|^p\,dx\) is weakly lower semi-continuous for
 \( n,p \) in Theorem \ref{main theorem 1} because
\begin{itemize}
       \item[(i)] if \(n\leq 4\), then  \(W^{2,2}(\Omega)\) is compactly imbedded
       into \( L^p(\Omega)\) for all \(  p>2\), and
       \item[(ii)] if \(5\leq n\leq 7\), then \(W^{2,2}(\Omega)\) is compactly imbedded
       into \( L^p(\Omega)\) for all \( 2<p< \frac{2n}{n-4}\).
\end{itemize}

It then follows for both of cases that
\[
\lim_{n\to \infty}\int_{\Omega}| u_m|^P=\int_{\Omega}| u|^P\,dx ,
\]

\item The map \( u\to\int_{\Omega}|u|^{p-2}uv\,dx\) is weakly lower semi-continuous for
\( n,p \) in Theorem \ref{main theorem 1} because
\begin{itemize}
   \item[(i)] if \(n\leq 4\), then  \(W^{2,2}(\Omega)\) is compactly  imbedded into
   \( L^{2(p-1)}(\Omega)\) for all \(  p>2\), and
   \item[(ii)] if \(5\leq n\leq 7\), then \(W^{2,2}(\Omega)\) is compactly imbedded
   into \( L^{2(p-1)}(\Omega)\) for all \( 2<p< \frac{2n-4}{n-4}\).
\end{itemize}
In both cases, we have  $|u|^{p-2}u\in L^2(\Omega) $ from which we deduce that the map
\( u\to\int_{\Omega}|u|^{p-2}uv\,dx\) is continuous functional and
\[
\lim_{n\to \infty}
\int_{\Omega}|u_m|^{p-2}uv\,dx= \int_{\Omega}|u|^{p-2}uv\,dx.
 \]
\end{itemize}
This completes the proof.
\end{proof}

We are now in a position to state the following result addressing condition (i)
in Theorem \ref{main theorem 1}.

\begin{lemma}\label{cond i}
Let \( K=K(r)\) be a convex and weakly closed subset of \(V\) defined in \eqref{def K}.
Let \( M : K \times K \to\mathbb{R}\) be defined as \eqref{def II} and \(n,p\) as in
Theorem \ref{main theorem 1}. Then there exists \(\bar{u} \in K\)    such that
\[M(\bar{u}, v)\leq 0\quad  \forall v\in  K.\]
\end{lemma}

\begin{proof}
 We shall show that the function \(M\) satisfies all the conditions of the Ky Fan's
Min-Max Principle presented in Lemma \ref{Ky Fan 2}.
The condition (1) is provided by Lemma \ref{WLSC}. For
each \(u\in K\), the map \(v \to M(u, v)\) is concave on \(K\) since \(M(u, v)\) is a
linear functional with respect to \(v\) except
\( -\frac{1}{2}\int_{\Omega} |\nabla v|^2\,dx\), which is in fact concave.
 Also we have \( M(u, u) = 0 = \gamma\) for every \(u \in K\). Finally, since
\(u\in K\) we have that \(\rVert u \rVert_{ W^{2,2}(\Omega)}\leq r\).
Thus, we can conclude that
\( \{u \in K : M(u, v) \leq 0\}\) is bounded. It now follows by Lemma \ref{Ky Fan 2}
that there exists \(\bar{u} \in K\) such that
\[
M(\bar{u}, v)\leq 0\quad  \forall v\in  K,
\]
as desired.
\end{proof}

Our next task consists of verifying condition (ii) in Theorem \ref{main theorem 1}.
To do this, we start with the following two lemmas, which provide us the required
estimates.
Hereafter $C$ will denote a positive constant, not necessarily the same one.

\begin{lemma}\label{C1C2}
 Let \(\Omega \subset \mathbb{R}^n \) be a bounded domain and \( 1<p\).
Then for any $\ u\in K(r) $ we have
\[
\big\rVert f+|u|^{p-1}u-\Lambda u\big\rVert_{L^2(\Omega)}
\leq \big( \rVert f\rVert_{L^2(\Omega)}
+\rVert u\rVert^{p-1}_{L^{2(p-1)}}+\rVert {u}\rVert^{2}_{W^{1,2^*}(\Omega)}
\rVert  {u}\rVert^2_{L^{\frac{n}{2}}(\Omega)} \big).
\]
\end{lemma}

\begin{proof}
Let \(u \in K(r)\). By H\"older's inequality we have
\begin{align*}
     \big\rVert f+|u|^{p-2}u-\Lambda u\big\rVert_{L^2(\Omega)}
     &\leq \rVert f\rVert_{L^2(\Omega)}+\rVert u^{p-1}\rVert_{L^2(\Omega)}
     +\rVert \Lambda u\rVert_{L^2(\Omega)}\\
     &\leq \rVert f\rVert_{L^2(\Omega)}+\rVert u\rVert^{p-1}_{L^{2(p-1)}(\Omega)}
     +\rVert \nabla u\rVert^2_{L^{2^*}(\Omega)}
     \rVert  u\rVert^2_{L^{\frac{n}{2}}(\Omega)}, \quad (\text{where }2^*=\frac{2n}{n-2})\\
     &\leq \rVert f\rVert_{L^2(\Omega)}+\rVert u\rVert^{p-1}_{L^{2(p-1)}(\Omega)}
     +\rVert  u\rVert^2_{W^{1,2^*}(\Omega)}\rVert  u\rVert^2_{L^{\frac{n}{2}}(\Omega)}.
\end{align*}
as desired.
\end{proof}

\begin{lemma}\label{rbar}
Let \( p>2 \) and \(C>0\) be given. Then there exists \(0<{r}\in \mathbb{R}\) which satisfies
 \[ C (\rVert f\rVert_{L^2(\Omega)} +{r}^{p-1}+{r}^4)\leq {r},\]
 where \(\rVert f\rVert_{L^2(\Omega)}\) be small enough.
\end{lemma}

\begin{proof} Since \(p >2\), we can choose \({r}\) such that
\[
C( {r}^{p-1}+ {r}^4)\leq \frac{{r}}{2}.
\]
Now if \(C\rVert f\rVert_{L^2(\Omega)}\leq \frac{{r}}{2}\) then we have
\[
C (\rVert f\rVert_{L^2(\Omega)} +{r}^{p-1}+{r}^4)\leq {r},
\]
as desired.
\end{proof}

Here is another useful results that we shall use in the sequel.
See \cite[Theorem 1.2]{RH} for a more general version of the following result.

\begin{lemma}\label{nabla lem}
If \(g\in L^2(\Omega)\), then there exists  \(u \in W^{2,2}(\Omega)\cap H_0^1(\Omega)\),
a scalar function \(P : \Omega \to\mathbb{R}\)  and a constant \( C\) such that
 \[
\Delta u + \nabla P = g, \quad  \nabla\cdot u=0,\quad  u|_{\partial \Omega}=0,
\]
and
\[\
rVert \Delta u \rVert_{L^2(\Omega)}+\rVert \nabla P \rVert_{L^2(\Omega)}
\leq C\rVert g \rVert_{L^2(\Omega)}.
\]
\end{lemma}

The following inequality is proved in \cite[Lemma 9.17]{DN}.

\begin{lemma} \label{Lu}
Let \(\Omega \) be a bounded \( C^{1,1 }\) domain in \(\mathbb{R}^n\) and let the
operator \(Lu = a^{ij}(x)D_{ij}u+b^i(x)D_iu+c(x)u\)
 be strictly Elliptic in \(\Omega \) with coefficients
 \( a_{ij} \in C(\Omega )\), \(b_i , c \in L^\infty(\Omega) \), with
 \(i, j = 1, \dots, n\) and \(c \leq 0\). Then there
 exists a positive constant \(C\) (independent of \(u\)) such that
\[
 rVert u \rVert_{ W^{2,p}(\Omega)}\leq C \rVert Lu \rVert_{ L^{p}(\Omega)},
\]
for all \(u \in  W^{2,p}(\Omega) \cap W_{0}^{1,p}(\Omega),\  1 < p < \infty\).
\end{lemma}

Here comes a direct consequence of Lemma \ref{Lu}.

\begin{corollary} \label{cor lap}
Let \(\Omega \) be a bounded \( C^{1,1 }\) domain in \(\mathbb{R}^n\).
Then there exists  a constant \(C\) such that
\[
\rVert u \rVert_{ W^{2,2}(\Omega)}\leq C \rVert \Delta u \rVert_{ L^2(\Omega)},
\]
for all \(u \in  W^{2,2}(\Omega) \cap H_0^1(\Omega)\).
\end{corollary}

We are now ready to prove Theorem \ref{main theorem 1}.

\begin{proof}[Proof of Theorem \ref{main theorem 1}]
 Without loss of generality we may suppose that \(\mu=-1\).
We define \(K = K({r})\) for $r>0$ to be determined presently.
By Lemma \ref{cond i} we have the existence of a non-trivial \(\bar{u} \in K\)
such that
\[
M(\bar{u}, v)\leq 0\quad  \forall v\in  K.
\]
Now we shall show the existence of \(\bar{v}\) that satisfy condition (ii) in the
theorem. Consider
\[
g(x)=f+|\bar{u}|^{p-1}\bar{u}-\Lambda \bar{u}.
\]

Thus we have to show there exists \(\bar{v}\in K\)  that the following equation holds
in the weak sense,
\begin{equation}\label{eq. lap}
   -\Delta v + \nabla P=g(x).
\end{equation}
By Lemma \ref{nabla lem} there exists \(\bar{v}\in V \)
which satisfies \eqref{eq. lap} and
\begin{equation}\label{nabla}
\rVert \Delta \bar{v} \rVert_{L^2(\Omega)}
+\rVert \nabla P \rVert_{L^2(\Omega)}\leq C\rVert g \rVert_{L^2(\Omega).}
\end{equation}
It is sufficient to show that \( \bar{v}\in K\). The  estimate \eqref{nabla}
together with  Lemma \ref{C1C2} imply that
\begin{equation} \label{nabla111}
\begin{aligned}
   \rVert \Delta {\bar{v}} \rVert_{L^2(\Omega)}
   +\rVert \nabla P \rVert_{L^2(\Omega)}
   &\leq C \rVert f+|\bar{u}|^{p-1}\bar{u}-\Lambda \bar{u}\rVert \\
   &\leq C\big( \rVert f\rVert_{L^2(\Omega)}+\rVert \bar{u}\rVert^{p-1}_{L^{2(p-1)}}
   +\rVert \bar{u}\rVert^{2}_{W^{1,2^*}(\Omega)} \rVert
    \bar{u}\rVert^2_{L^{\frac{n}{2}}(\Omega)} \big) .
\end{aligned}
\end{equation}

On the other hand, Corollary \ref{cor lap} together with \eqref{nabla111} yield that
\begin{equation}\label{CC}
\begin{aligned}
    \rVert \bar{v} \rVert_{ W^{2,2}(\Omega)}
     &\leq C \rVert \Delta \bar{v} \rVert_{ L^2(\Omega)}
     \leq {C}\big(\rVert \Delta \bar{v} \rVert_{L^2(\Omega)}
    +\rVert \nabla P \rVert_{L^2(\Omega)} \\
     &\leq C\big( \rVert f\rVert_{L^2(\Omega)}
    +\rVert \bar{u}\rVert^{p-1}_{L^{2(p-1)}}
    +\rVert \bar{u}\rVert^{2}_{W^{1,2^*}(\Omega)} \rVert
    \bar{u}\rVert^2_{L^{\frac{n}{2}}(\Omega)} \big).
\end{aligned}
\end{equation}


From the  imbeddings of  \(W^{2,2}(\Omega)\hookrightarrow L^{2(p-1)}\) and
\(W^{2,2}(\Omega)\hookrightarrow W^{1,2^*}(\Omega)\)  we obtain from \eqref{CC}
that
\begin{equation}\label{CC2}
  \rVert \bar{v} \rVert_{ W^{2,2}(\Omega)}
  \leq C\big( \rVert f\rVert_{L^2(\Omega)}+\rVert \bar{u}\rVert^{p-1}_{w^{2,2}(\Omega)}
  +\rVert \bar{u}\rVert^{2}_{W^{2,2}(\Omega)} \rVert  \bar{u}\rVert^2_{W^{2,2}(\Omega)}
  \big).
\end{equation}
Let $r$ be as in Lemma \ref{rbar} for $C$ given in the last inequality above.
The inequality \eqref{CC2} and Lemmas \ref{rbar} yield that
\[
  \rVert \bar{v} \rVert_{ W^{2,2}(\Omega)}
\leq   C(\rVert f\rVert_{L^2(\Omega)} +{r}^{p-1}+{r}^4)\leq {r},
\]
where \(\rVert f\rVert_{L^2(\Omega)}\) is small enough.
That means \( \bar{v}\in K\) and so $\bar{v}=\bar{u}$.
This completes the proof of (i) and (ii).
Now  the inequality \eqref{nabla111}  concludes that
  \begin{equation} \label{nabla11}
  \rVert \Delta {u} \rVert_{L^2(\Omega)}+\rVert \nabla P \rVert_{L^2(\Omega)}
  \leq C\big( \rVert f\rVert_{L^2(\Omega)}+\rVert u\rVert^{p-1}_{L^{2(p-1)}}
  +\rVert {u}\rVert^{2}_{W^{1,2^*}(\Omega)}
  \rVert  {u}\rVert^2_{L^{\frac{n}{2}}(\Omega)} \big) .
\end{equation}
\end{proof}

\section{Proof of Theorems \ref{main theorem 2} and \ref{main theorem 3}}

We need some preliminary results before proving  the theorems in this section.
We consider  the same notation  for the Banach space
\( V= \{ u\in H_0^1(\Omega)\cap L^p(\Omega),\nabla . u=0\} \)
with norm \( \rVert u \rVert= \rVert u \rVert_{ H_0^{1}(\Omega)}
+\rVert u \rVert_{ L^{p}(\Omega)}\).
Where $\Omega$ is a  bounded domain of $\mathbb{R}^n$, the operator
\(\Lambda u = (u . \nabla)u\) may  not be defined on whole space $H_0^1(\Omega)$.
Although, there exist constant $C$ such  that
\[
 |\langle\Lambda u, v\rangle|  =\Big|\int_\Omega\sum_{j,k=1}^n u_k
 \frac{\partial u_j}{\partial x_k}v_j\Big|
 \leq  C\rVert u \rVert_{L^2(\Omega)}\rVert \nabla u \rVert_{L^2(\Omega)}
  \rVert v \rVert_{C^1(\Omega)}\,.
\]
which means that for the dense linear subspace
\[
E = \{ u\in C_c^1(\Omega),\nabla .  u=0\}
\]
of $V$, we have that $\Lambda$ is well defined.
We shall define \(\Phi: V\to\mathbb{R}\)  by
\[
\Phi(u)= \frac{1}{2}\int_{\Omega} |\nabla u|^2\,dx+\frac{1}{p}\int_\Omega
|u|^{p}\,dx -\int_{\Omega} f\;u\,dx.
\]
We also define \(H:  V \times  V \to\mathbb{R}\)  by
\[
H(v, u) = \Phi(u) -\Phi(v).
\]
For \(r > 1\), we set
\[
K(r) = \{u\in V; \rVert u\rVert \leq  r\},
\]
that is convex and  weakly closed in \(V\) by similar arguments as in the proof
of Lemma \ref{Kclosed}. Let
\[
K_0(r) = K(r) \cap E,
\]
and define \(M : K(r) \times K_0(r) \to\mathbb{R}\) by
\begin{equation}\label{defM2}
  M(u, v) = H(v, u) -\langle \Lambda u,v \rangle.
\end{equation}

When $\mu>0$ in the damping term,
we shall  use a different
 version of Ky-Fan minimax theorem
(See \cite[Lemma 12.1]{NG}) for a proof).
This version  is more practical when one expects less regularity of  the solution.
For a subset set $D$,  we denote its convex hull by $conv(D)$.

\begin{lemma} \label{Ky Fan 1}
Let \(\emptyset \neq D\subset E \subset H\) where \(E\) is  a weakly compact convex
set in a Banach space \(H\), and consider
\(M : E \times conv(D)\to{\mathbb{R}}\) to be a function such that:
\begin{itemize}
    \item[(1)] For each \(y\in D\), the map \(x \to M(x, y)\) is weakly lower semi-
    continuous on \(E\).
    \item[(2)]  For each \(x \in E\), the map \(y \to M(x, y)\) is concave on \(conv(D)\).
    \item[(3)]  \(M(x, x) \leq 0\) for every \(x \in conv(D)\).
\end{itemize}
Then there exits \(\bar{x} \in E\) such that \(M(\bar{x}, y) \leq 0\) for all
\(y \in D\).
\end{lemma}

\begin{proof}[Proof of Theorem \ref{main theorem 2}]
Without loss of generality we may suppose that \(\mu=1\). By  similar arguments as
in Lemma \ref{WLSC} for \(M\) defined in \eqref{defM2} we obtain that
\begin{itemize}
   \item   For each \(v \in K_0(r)\) the function \(u \to M(v, u)\) is weakly lower semi-continuous.
   \item  For each \(u \in K(r)\) the function \(v \to M(v, u)\) is concave.
   \item \(M(u,u)=0, \forall u\in K_0(r)\)
\end{itemize}
Now we can apply  Ky-Fan minimax
principle (Lemma \ref{Ky Fan 1}), which yields there exits \( \bar{u}_r\in K(r)\)
such that
\begin{equation}\label{ine. M}
  M(\bar{u}_r,v)=H(v, \bar{u}_r) -\langle\Lambda \bar{u}_r,v \rangle
  \leq 0,\quad \;\forall v\in K_0(r)
\end{equation}
Substituting \(v =0\)  in the latter inequality  implies that \(\Phi(\bar{u_r})\leq0\).
Now the coercivity of the functional \(\Phi\)  follows that
\(\{\bar{u}_r\}_r\) is bounded in \(V\) and so  there
exists a sequence \(r_n \to\infty\) and \(\bar{u}\in V\)  such that
\(\bar{u}_{r_n}\to\bar{u}\) weakly in \(V\). If \(v\in E\) is fixed, then
from \eqref{ine. M} and the weak lower semi-continuity of the functions involved, we obtain
\begin{equation}\label{ine. H}
  H(v, \bar{u}) -\langle\Lambda \bar{u},v \rangle\leq 0
\end{equation}
This indeed implies that
\begin{equation}\label{ine. sup 2}
  \sup_{v\in E,\rVert v\rVert_{E}\leq 1}
  \langle \Lambda \bar{u},v\rangle +\inf_{\rVert z \rVert\leq 1} H(z,\bar{u}) \leq 0
\end{equation}
Therefore,
\begin{equation}
  \sup_{v\in E,\rVert v\rVert_{E}\leq 1}   \langle \Lambda \bar{u},v\rangle
  \leq -\inf_{\rVert z \rVert\leq 1} H(z,\bar{u}) \leq \infty.
\end{equation}
This implies that the linear functional \(l : E \to\mathbb{R}\) defined by
\(l(v) = <\Lambda \bar{u}, v>\) is continuous. It now follows from the bounded
linear extension theorem that \(l\) can be extended to a bounded linear operator
\(L : V \to\mathbb{R}\) with the same
operator norm as \(l\). It then follows that there exists
\(\hat{\Lambda}\bar{u}\in V^*\) such that
\begin{equation}\label{hat}
\langle \hat{\Lambda} \bar{u},v\rangle=<\Lambda \bar{u},v>,\quad  \forall v\in E.
\end{equation}
This together with \eqref{ine. H} yield that
\begin{equation}\label{hat 2}
H(v, \bar{u}) -\langle \hat{\Lambda} \bar{u},v\rangle \leq 0 ,\quad  \forall v\in E.
\end{equation}
But since \(E\) is dense in \(V\) and  expression \eqref{hat 2} is continuous with
respect to \(v\), we can conclude that
\begin{equation}\label{hat 3}
H(v, \bar{u}) -\langle \hat{\Lambda} \bar{u},v\rangle \leq 0 ,\quad  \forall v\in V.
\end{equation}
Now by substituting \(v=\bar{u}+t\eta,\;\eta \in V\), into \eqref{hat 3}  we obtain that
\begin{equation}\label{hat 4}
H(\bar{u}+t\eta, \bar{u}) -\langle \hat{\Lambda} \bar{u},\bar{u}+t\eta\rangle
\leq 0 ,\quad  \forall t \in \mathbb{R}.
\end{equation}
Dividing \eqref{hat 4} by  \(t>0\) and letting \(t\) converge to zero yields that
\begin{equation} \label{hat 5}
  \int_\Omega\nabla \bar{u}.\nabla\eta\,dx + \int_\Omega |\bar{u}|^{p-2}\bar{u} \eta \,dx
   -\int_\Omega f\eta\,dx + \int_\Omega \hat{\Lambda} \bar{u}\;\eta \geq 0,\quad
   \forall\eta\in  V.
\end{equation}
Now substituting \(\eta\) by \(-\eta\)  in \eqref{hat 5} we deduce the opposite
inequality and thus
\begin{equation} \label{hat 6}
  \int_\Omega\nabla \bar{u}.\nabla\eta\,dx
  + \int_\Omega |\bar{u}|^{p-2}\bar{u} \eta \,dx
  -\int_\Omega f\eta\,dx + \int_\Omega \hat{\Lambda} \bar{u}\;\eta = 0,
  \quad \forall\eta\in  V.
\end{equation}
This together with \eqref{hat} follow that
\begin{align*} %\label{hat 7}
  \int_\Omega\nabla \bar{u}.\nabla\eta\,dx
  + \int_\Omega |\bar{u}|^{p-2}\bar{u}\; \eta \,dx
  -\int_\Omega f\eta\,dx + \int_\Omega {\Lambda} \bar{u}\;\eta = 0,\quad
  \forall\eta\in  E.
\end{align*}
This completes the proof of part (i).

For the proof of the second part  we consider two cases \(n\leq 4\) and \(p\geq 4\)
separately.
\smallskip

\noindent\textbf{Case 1: \((n\leq 4)\).} Let \(v\in V\). If \( n \leq 4\), then
\( 4 \leq\frac{2n}{n-2}\). From  continuous imbedding of Sobolev space
\(H_0^1(\Omega)\)   into \(L^4(\Omega)\), and by the H\"{o}lder inequality
we obtain for operator \(\Lambda u\)
\begin{align*}
  |\langle\Lambda u, v\rangle|
  =\big|\int_\Omega\sum_{j,k=1}^n u_k\frac{\partial u_j}{\partial x_k}v_j\big|
  &\leq  C\rVert u v\rVert_{L^2(\Omega)}\rVert \nabla u \rVert_{L^2(\Omega)}\\
  &\leq  C\rVert u \rVert_{L^4(\Omega)}\rVert v\rVert_{L^4(\Omega)}
  \rVert  u \rVert_{H_0^1(\Omega)}< \infty.
\end{align*}
This means, the operator \(\Lambda u\) is well defined on \(V\).
Since \( E\) is a dense subspace of \(V \), from uniqueness of the bounded linear
extension theorem we have \(
\langle \hat{\Lambda} \bar{u},v\rangle=<\Lambda \bar{u},v>,\; \forall v\in V.
\)
Now the result follows from \eqref{hat 6}.
\smallskip

\noindent\textbf{Case 2: \(p\geq 4\).}
 For \(v\in V\), since \(V\subset L^p(\Omega)\)  we can deduce that
\[
  |\langle\Lambda u, v\rangle|
  \leq  C\rVert u v\rVert_{L^2(\Omega)}\rVert \nabla u \rVert_{L^2(\Omega)}
  \leq  C\rVert u \rVert_{H_0^1(\Omega)} \rVert u \rVert_{L^p(\Omega)}
  \rVert v \rVert_{L^{\frac{2p}{p-2}}(\Omega)}<\infty
\]
where the last inequality follows from
\( \frac{2p}{p-2}\leq p\). Thus operator \(\Lambda u\) is well defined on whole
\(V\) and this completes the proof.
\end{proof}

We would like to remark that in the last part of the proof we make extensive use of \( \frac{2p}{p-2}\leq p\). Note that if \(2<p<4\) and \(n\geq 5\), then \(\frac{2p}{p-2}>p\), and there is no guarantee for
\(\rVert v\rVert_{L^{\frac{2p}{p-2}}(\Omega)}<\infty\).\\

As we have just seen, the case of $\mu>0$ with the linear damping term was covered
in the  theorem \ref{main theorem 2}. But for $\mu<0$, due to  an essential role of
Lemma \ref{rbar} in the proof of theorem \ref{main theorem 1} we were not  be able to
deal with the linear damping term in this theorem. However, with a similar argument to
theorem \ref{main theorem 2}, we would be in a position to manage it separately.
Note that when $p=2$ we have that \( V = \{ u\in H_0^1(\Omega),\nabla u=0\} \), and
\[
\Phi(u)= \frac{1}{2}\int_{\Omega} |\nabla u|^2\,dx
+\frac{\mu}{2}\int_\Omega   |u|^{2}\,dx -\int_{\Omega} f\;u\,dx.
\]

\begin{proof}[Proof of Theorem \ref{main theorem 3}]
In the same way as in proof of Theorem \ref{main theorem 2}, it follows from  the
Ky-Fan minimax principle (Lemma \ref{Ky Fan 1}) that there exits
\( \bar{u}_r\in K(r)\) with
\begin{equation}\label{ine. M-2}
 M(\bar{u}_r,v)=H(v, \bar{u}_r) -\langle\Lambda \bar{u}_r,v \rangle \leq 0,\quad
 \forall v\in K_0(r).
\end{equation}
Now we claim that \(\{\bar{u}_r\}_r\) is bounded in \(V\) and so there exists a sequence
\(r_n \to\infty\) and \(\bar{u}\in V\)  such that \(\bar{u}_{r_n}\rightharpoonup \bar{u}\) weakly in \(V\).
Thus, by similar arguments as in proof of Theorem \ref{main theorem 2} we obtain the
result.

Now  to complete the proof we have to show the claim.
Assume, by contradiction, that \(\{\bar{u}_r\}_r\) is unbounded. So  there exists a
sequence \(r_m \to\infty\)  such that \(\{\bar{u}_{r_m}\}_m\) is unbounded.
By substituting \(v=0\)  in  \eqref{ine. M-2} we obtain that
\begin{equation}\label{claim 1}
 \Phi(\bar{u}_{r_m})=  \frac{1}{2}\int_{\Omega} |\nabla \bar{u}_{r_m}|^2\,dx
 +\frac{\mu}{2}\int_\Omega   |\bar{u}_{r_m}|^{2}\,dx
 -\int_{\Omega} f\;\bar{u}_{r_m}\,dx\leq0.
\end{equation}

Let \( t^2_m=\int_{\Omega}|\nabla \bar{u}_{r_m}|^2\,dx\), and
\( w_m=\frac{\bar{u}_{r_m}}{t_m}\).
Note that \(\rVert w_m\rVert_{H_0^1(\Omega)} =1\). Thus, there exists a
\(w=(z_1,\dots,z_n)\in V\) such that \(w_m \rightharpoonup w\) weakly in \(V\).
It follows that \(w \neq 0\), because  dividing \eqref{claim 1} by  \(t^2_m\)
we obtain
\begin{equation}\label{claim 2}
 \frac{1}{2}+\frac{\mu}{2}\int_\Omega   |w_m|^{2}\,dx
 \leq \frac{1}{t_m}\int_{\Omega} f\;w_m\,dx,
\end{equation}
and letting \(m\to \infty\), due to the compact imbedding
\(H_0^1(\Omega)\hookrightarrow L^2(\Omega)\) we obtain that
\begin{equation}\label{claim 3}
 \frac{1}{2}+\frac{\mu}{2}\int_\Omega   |w|^{2}\,dx \leq 0,
\end{equation}
which implies \(w\neq 0\). Also, we have
\begin{align*}
  \frac{1}{2}\int_{\Omega} |\nabla w|^2\,dx+\frac{\mu}{2}\int_\Omega  |w|^{2}\,dx
  &\leq \liminf_{m\to \infty}\big(\frac{1}{2}\int_{\Omega} |\nabla w_m|^2\,dx+\frac{\mu}{2}\int_\Omega  |w_m|^{2}\,dx\big)\\
  &\leq \frac{1}{2}+\frac{\mu}{2}\int_\Omega   |w|^{2}\,dx.
\end{align*}
This estimate  together with  \eqref{claim 3} yield that
\[
  \int_{\Omega} |\nabla w|^2\,dx+{\mu}\int_\Omega  |w|^{2}\,dx \leq 0.
\]
Therefore,
\begin{equation}\label{claim 4}
  \frac{\int_{\Omega} |\nabla w|^2\,dx}{\int_\Omega  |w|^{2} \,dx} \leq -\mu,
\end{equation}
from which  with hypothesis \(-\lambda_1\leq \mu\) of theorem we obtain that
\begin{equation}\label{claim 5}
  \frac{\int_{\Omega} |\nabla w|^2\,dx}{\int_\Omega  |w|^{2} \,dx} \leq \lambda_1.
\end{equation}
On the other hand, for  first eigenvalue \(\lambda_1\) of \(-\Delta\) we have
\begin{equation}\label{claim 6}
  \frac{\int_{\Omega} |\nabla w|^2\,dx}{\int_\Omega  |w|^{2}\,dx}= \frac{\sum_{i=1}^n\int_{\Omega} |\nabla z_i|^2\,dx}{\sum_{i=1}^n\int_\Omega  z_i^{2}\,dx} \geq \frac{\sum_{i=1}^n \lambda_1 \int_\Omega  z_i^{2}\,dx}{\sum_{i=1}^n\int_\Omega  z_i^{2} \,dx}=\lambda_1,
\end{equation}
where $w=(z_1,\dots,z_n)$.
It then follows from \eqref{claim 5}  and \eqref{claim 6} that
\begin{equation}\label{claim 7}
   \lambda_1 =  \frac{\int_{\Omega} |\nabla w|^2\,dx}{\int_\Omega  |w|^{2}\,dx},
\end{equation}
from which we obtain that
\[\int_{\Omega} |\nabla z_i|^2\,dx=\lambda_1{\int_\Omega  z_i^{2}\,dx},  \quad (i=1,..,n).\]
Therefore,
\begin{equation}\label{claim 8}
   -\Delta z_i  = \lambda_1 z_i,\quad i=1,\dots,n.
\end{equation}
Since the first eigenvalue of the $-\Delta$ is simple it  follows  that there exists
\(\alpha=(\alpha_1,\dots,\alpha_n) \in \mathbb{R}^n\) such that
\begin{equation}\label{claim 9}
    z_i  = \alpha_i \psi_1,\quad i=1,\dots,n,
\end{equation}
where \(\psi_1>0\) is the  unique eigenfunction of $-\Delta$ corresponding to
\(\lambda_1\) with \(\rVert\psi_1\rVert_{L^2(\Omega)}=1\), i.e.
\[
-\Delta \psi_1  = \lambda_1 \psi_1, \quad  \psi_1 |_{\partial \Omega}=0.
\]
Since \(\nabla.w=0\), it follows from \eqref{claim 9} that
\begin{equation}\label{claim 10}
   0=\sum_{i=1}^n \frac{\partial z_i}{\partial x_i}= \sum_{i=1}^n\alpha_i\frac{\partial \psi_1}{\partial x_i}=\alpha .\nabla\psi_1.
\end{equation}

Now let \(x\)  be an interior point of \(\Omega\) and \( \bar{x} \) the closest point
on \(\partial\Omega\) to \(x\) such that  \(\bar{x}-x=C\alpha\) for some constant
\(C\in \mathbb{R}\), and the line joining $x$ to $\bar x$ lies in $\bar \Omega$. Define
\( g : [0,1]\to\mathbb{R}\) by
\[
 g(t)=\psi_1(tx+(1-t)\bar{x}).
\]
It can be easily deduced from \eqref{claim 10} that
\[
  g'(t)=(x-\bar{x}).\nabla\psi_1(tx+(1-t)\bar{x})
  =C\alpha.\nabla\psi_1(tx+(1-t)\bar{x})=0.
\]
Thus, \(g\) is a constant function and since \(\psi_1 |_{\partial \Omega}=0\) we have
\[
g(t)= g(0)=\psi_1(\bar{x})=0,  \quad \forall t\in [0,1],
\]
which implies that \(\psi_1(x)=0 \). This is the contradiction we wanted.
\end{proof}

\subsection*{Acknowledgments}
A. Moameni acknowledges support of the Natural Sciences and Engineering Research 
Council of Canada.
A. Khatib and S. Mousavinasr are grateful to CNPq (401772/2022-5) and 
Coordena\c{c}\~ao de Aperfei\c{c}oamento de Pessoal de N\'ivel Superior (CAPES), 
Brazil, for their support.


\begin{thebibliography}{10}


\bibitem{HB}
H. Brezis, L. Nirenberg, G. Stampacchia;
\emph{A remark on Ky Fan's Minimax Principle},  Bollettino U. M. I
(1972), 293-300

\bibitem{XC}
X. Cai and Q. Jiu;
\emph{Weak and strong solutions for the incompressible Navier-Stokes equations
with damping}, J. Math. Anal. Appl., 343 (2008), 799-809.

\bibitem{RH}
R. Farwig and H. Sohr;
\emph{Generalized resolvent estimates for the Stokes system in bounded and unbounded
domains}, J. Math. Soc. Japan, vol 46 no 4 (1994), 607-643.

\bibitem{DN} D. Gilbarg, N. S. Trudinger;
\emph{Elliptic partial differential equations of second order}, Reprint of the (1998)
edition. Classics in Mathematics. Springer-Verlag, Berlin, (2001).

\bibitem{NG} N. Ghoussoub;
\emph{Self-dual Partial Differential Systems and Their Variational Principles},
Springer Monographs in Mathematics, Springer, New York (2008).

\bibitem{EH}
E. Hopf;
\emph{Über die Anfangswertaufgabe für die hydrodynamischen Grundgleichungen},
 Math. Nachr. 4 (1951), 213-231.

\bibitem{LHS}
L. Hsiao;
emph{Quasilinear Hyperbolic Systems and Dissipative Mechanisms}, World Scientific, 1997.

\bibitem{YJX}
Y. Jia, X. W. Zhang, B. Q. Dong;
\emph{The asymptotic behavior of solutions to three dimensional Navier-Stokes equations
with nonlinear damping}, Nonlinear Anal. Real World
Appl., 12 (2011), 1736-1747.

\bibitem{ZHJ}
Z. H. Jiang;
\emph{Asymptotic behavior of strong solutions to the 3D Navier-Stokes equations with
a nonlinear damping term}, Nonlinear Anal., 75 (2012), 5002-5009.

\bibitem{ZHM}
Z. H. Jiang, M. X. Zhu;
\emph{The large time behavior of solutions to 3D Navier-Stokes equations
with nonlinear damping}, Math. Methods Appl. Sci., 35 (2012), 97-102.


\bibitem{JL}
J. Leray;
\emph{Sur le mouvement d'un liquide visqueux emplissant l'espace},
 Acta Math. 63 (1934), 193-248.

\bibitem{ZLD}
Z. Li, D. Shi, M. Li;
\emph{Stabilized mixed finite element methods for the Navier-Stokes equations with
damping}, Math. Methods Appl. Sci. 42 (2019), 605-619.

\bibitem{WLW}
W. Li, W. Wang, Q. Jiu;
\emph{Existence and uniqueness of the weak solutions for the steady
incompressible Navier Stokes equations with damping}, African diaspora journal of
mathematics, volume 12, number 2 (2011), 57-72.

\bibitem{LML}
M. Li, Z. Li, D. SHI;
\emph{Unconditional Optimal Error Estimates for the Transient Navier-Stokes Equations
with Damping}. Adv. Appl. Math. Mech., v. 14, n. 1 (2022), 248-274.

\bibitem{HLH}
H. Liu and H.J. Gao;
\emph{Decay of solutions for the 3D Navier-Stokes equations with damping}, Appl.
Math. Lett., 68 (2017), 48-54.

\bibitem{AMO}
 A. Moameni;
\emph {Critical point theory on convex subsets with applications in differential
 equations and analysis.} J. Math. Pures Appl. (9) 141 (2020), 266-315.
	

\bibitem{HBO}
H. B. de Oliveira;
\emph{Existence of weak solutions for the generalized Navier-Stokes equations with
damping}, Nonlinear Diff. Eqs. Appl., 20 (2013) 797-824.

\bibitem{HQL}
H. Qiu, L. Mei;
\emph{Multi-level stabilized algorithms for the stationary incompressible Navier-Stokes
equations with damping}, Appl. Numer. Math., 143 (2019), pp.188-202.

\bibitem{HQY} H. Qiu, Y. Zhang, L. Mei;
 \emph{A mixed-FEM for Navier-Stokes type variational
inequality with nonlinear damping term}, Comput. Math. Appl., 73(10) (2017),
2191-2207

\bibitem{XLS}
X. L. Song, F. Liang, J. H. Wu;
\emph{Pullback D-attractors for three-dimensional Navier-Stokes
equations with nonlinear damping}, Bound. Value Probl., (2016), Paper No. 145, 15 pp.

\bibitem{XLY}
X.L. Songm Y. R. Hou;
\emph{Uniform attractor for three-dimensional Navier-Stokes equations with nonlinear
damping}, J. Math. Anal. Appl., 422 (2015), 337-351.

\bibitem{XSY}
X. Song, Y. Hou;
\emph{Attractors for the three-dimensional incompressible Navier-Stokes equations
with damping}, Discrete Contin. Dyn. Syst. Ser. A, 31 (2011), 239-252.

\bibitem{RTE}
R. Temam;
\emph{Navier-Stokes Equations, Theory and Numerical Analysis.
North-Holland Publishing Company}, Amsterdam, New York, Oxford, 1979.

\bibitem{YZ}
Y. Zhou;
\emph{Regularity and uniqueness for the 3D incompressible Navier-Stokes equations with
damping}, Appl. Math. Lett., 25 (2012), 1822-1825.

\bibitem{ZJZ}
Z. J. Zhang, X. L. Wu, M. Lu;
\emph{On the uniqueness of strong solution to the incompressible
Navier-Stokes equations with damping}, J. Math. Anal. Appl., 377 (2011), 414-419.

\bibitem{ZZC}
Z. Zhang, C. P. Wu, Z. Yao;
\emph{Remarks on global regularity for the 3D MHD system
with damping}, Appl. Math and Com., 333: 1-7 (2018)

\bibitem{ZBS}
B. Zheng, Y. Shang;
\emph{Two-level defect-correction stabilized algorithms for the simulation of
2D/3D steady Navier-Stokes equations with damping}.
Applied Numerical Mathematics, v. 163 (2021) 182-203.

\end{thebibliography}

\end{document}

